\documentclass[11pt]{letter}
\usepackage[margin=1in]{geometry}
\usepackage[hidelinks]{hyperref}
\signature{Paidi Geetesh Chandra\\(on behalf of all authors)}
\address{Paidi Geetesh Chandra\\Independent Researcher, India\\paidigeeteshchandra@gmail.com}
\begin{document}
\begin{letter}{The Editor-in-Chief\\Advances in Space Research}
\opening{Dear Editor,}

We submit the manuscript \emph{``Global Mg/Si and Al/Si maps of the Moon from Chandrayaan-2 CLASS with time-resolved solar excitation from GOES XRS''} by S.~Kanjilal, S.~Maji and P.~G.~Chandra, for consideration as a research article in \emph{Advances in Space Research}.

The Chandrayaan-2 CLASS Level-1 archive is public, but turning its spectra into composition requires the incident solar X-ray spectrum. That spectrum is often approximated or ignored on the assumption that element/Si ratios cancel it. Our paper makes three contributions that we believe will interest ASR readers working on planetary X-ray fluorescence and on Chandrayaan-2 data:
\begin{itemize}
\item We fit 1.05 million CLASS dayside spectra with a Poisson likelihood and a simplified fundamental-parameter model, using a time-resolved CHIANTI solar spectrum constrained by the GOES XRS temperature. The whole chain is reproducible from public data. We show quantitatively that solar line emission dominates the excitation of Mg, Al and Si, so that ratios do not cancel the solar spectrum.
\item We identify a solar-driven instrumental Al K$\alpha$ component that appears only when the Sun illuminates the spacecraft from behind. We show that it must be kept out of the dayside background model, and that doing so raises the agreement of Al/Si with Lunar Prospector GRS from $r=0.35$ to $0.72$.
\item From synthetic recovery tests we derive an explicit error budget that limits what can be recovered without a measured solar spectrum to Mg/Si and Al/Si. We present global 2\textdegree{} maps of these ratios with quality layers. The Al/Si map resolves individual maria and is validated against both GRS ($r=0.72$) and nine landing-site soils ($r=0.81$). Mg/Si agrees with GRS ($r=0.61$), and we report candidly that its landing-site test is inconclusive, as it is for GRS itself.
\end{itemize}

All analysis choices were fixed before validation. The history of the final background model, which was motivated by the diagnosis of an earlier model's validation but specified before it was run, is disclosed in full in the text, together with the earlier results. Code and per-spectrum results are public.

An earlier version of this work, which used a different and less complete physical model together with an image-fusion super-resolution step, was submitted to \emph{IEEE Transactions on Geoscience and Remote Sensing} and was not accepted. The present manuscript is a substantial rework that responds to that review: the forward model, the statistics, the background treatment and the validation have been rebuilt, and the super-resolution component has been removed. The manuscript is not under consideration elsewhere, and all authors have approved the submission. The authors declare no competing interests.

Possible reviewers with relevant expertise include specialists in planetary X-ray fluorescence (C1XS, CLASS, MESSENGER XRS) and in solar soft X-ray spectroscopy.

\closing{Sincerely,}
\end{letter}
\end{document}
